\chapter{Rank one, Casimir, and complete reducibility}
\label{ch:rank-one}

\section{The rank-one calculation}

Write \(\mathfrak{sl}_2(k)\) with basis
\[
 e=\begin{pmatrix}0&1\\0&0\end{pmatrix},\qquad
 f=\begin{pmatrix}0&0\\1&0\end{pmatrix},\qquad
 h=\begin{pmatrix}1&0\\0&-1\end{pmatrix}.
\]
Thus \([h,e]=2e,[h,f]=-2f,[e,f]=h\). In characteristic
zero this algebra is simple: an ideal is stable under
\(\operatorname{ad}h\), whose three distinct eigenspaces
are \(ke,kf,kh\). Polynomial projections onto these
eigenspaces show that a nonzero ideal contains one of
the basis vectors. Its brackets with the others then
give all three.

\begin{theorem}[Simple \(\mathfrak{sl}_2\)-modules]\label{res:C01}
Over an algebraically closed characteristic-zero field
the nonzero finite-dimensional simple
\(\mathfrak{sl}_2\)-modules are exactly \(L_n\),
\(n\in\mathbb Z_{\geq0}\). They have basis
\(v_0,\ldots,v_n\) on which
\[
 hv_j=(n-2j)v_j,\quad
 fv_j=v_{j+1},\quad
 ev_j=j(n-j+1)v_{j-1},
 \tag{3.1}\label{eq:sl2-basis}
\]
where \(v_{-1}=v_{n+1}=0\). They can be realized on
the degree-\(n\) homogeneous polynomials in \(X,Y\) by
\[
 e=X\partial_Y,\qquad f=Y\partial_X,\qquad
 h=X\partial_X-Y\partial_Y.
\]
\end{theorem}

\begin{proof}
Choose an \(h\)-eigenvector in a simple module. The
relation with \(e\) shows that every nonzero successive
\(e\)-image has its eigenvalue increased by \(2\).
There are only finitely many eigenvalues, so the last
nonzero image is a vector \(v\) with
\(ev=0,hv=\lambda v\).
Commuting \(h\) and \(e\) past powers of \(f\) gives
\[
 hf^jv=(\lambda-2j)f^jv,\qquad
 ef^jv=j(\lambda-j+1)f^{j-1}v.
 \tag{3.2}\label{eq:sl2-ladder}
\]
For the second formula, use
\([e,f^j]=\sum_{a=0}^{j-1}f^a h f^{j-1-a}\);
on \(v\) its coefficient is
\(\sum_{a=0}^{j-1}(\lambda-2(j-1-a))
=j(\lambda-j+1)\).

PBW with the order \(f,h,e\) says that the vectors
\(f^jv\) span the generated submodule, hence the whole
simple module. Their distinct \(h\)-eigenvalues make
the nonzero ones independent. Let \(n\) be the final
index with \(f^nv\ne0\). Applying \(e\) to
\(f^{n+1}v=0\) gives
\((n+1)(\lambda-n)f^nv=0\), so \(\lambda=n\).
This proves \eqref{eq:sl2-basis} and uniqueness.

Conversely the three polynomial operators satisfy
the displayed brackets. Set
\[
 v_j=f^jX^n=\frac{n!}{(n-j)!}X^{n-j}Y^j.
\]
Direct differentiation gives
\eqref{eq:sl2-basis}. A nonzero submodule contains an
\(h\)-eigenvector: the distinct-eigenvalue projections
are polynomials in \(h\). Each eigenline is \(kv_j\);
applying \(e^j\), whose coefficient is nonzero,
gives \(v_0\), and applying \(f\) gives every \(v_j\).
Thus the module is simple. This argument includes
\(n=0\), the trivial one-dimensional module.
\end{proof}

\section{Central operators and cohomology}

We now return to an arbitrary characteristic-zero field
and a finite-dimensional semisimple Lie algebra
\(\mathfrak g\). Its modules in this section are
finite dimensional. No root decomposition or highest
weight classification is needed for complete
reducibility.

\begin{theorem}[The module-dependent Casimir]\label{res:C02}
Let \(M\) be a nontrivial simple \(\mathfrak g\)-module.
There is a central element \(C_M\in U(\mathfrak g)\)
which annihilates the trivial module and acts
invertibly on \(M\).

More explicitly, let
\(\mathfrak h=\rho(\mathfrak g)\subseteq\operatorname{End}_k(M)\).
The trace form
\(B_M(x,y)=\operatorname{tr}_M(xy)\) is nondegenerate
on \(\mathfrak h\). Lift \(\mathfrak h\) to the
complementary ideal of \(\ker\rho\) in \(\mathfrak g\).
For a basis \(a_i\) and its \(B_M\)-dual basis \(a^i\)
of that ideal, one may take
\[
 C_M=\sum_i a_i a^i.
 \tag{3.3}\label{eq:module-casimir}
\]
Its trace on \(M\) is \(\dim\mathfrak h\).
If \(k\) is algebraically closed, it acts as the
nonzero scalar \(\dim\mathfrak h/\dim M\).
\end{theorem}

\begin{proof}
The image \(\mathfrak h\) is semisimple by
Theorem~\ref{res:B06}. The radical \(I\) of its
invariant trace form is an ideal. On \(I\), the
trace pairing in the faithful representation on
\(M\) is zero. In particular it vanishes on
\([I,I]\times I\), so the matrix trace criterion
shows that \(I\) is solvable. Semisimplicity of
\(\mathfrak h\) gives \(I=0\).

The element \(\sum_i a_i\otimes a^i\) is independent
of the dual basis: under the identification of one
factor with its dual by \(B_M\), it corresponds to
the identity endomorphism. Invariance of \(B_M\)
makes this tensor invariant under the diagonal
adjoint action. Concretely, if
\([x,a_i]=\sum_j b_{ji}a_j\), then
\([x,a^j]=-\sum_i b_{ji}a^i\), and the two sums
in its derivative cancel. Multiplication of the
tensor gives
\([x,C_M]=0\).
The other ideal \(\ker\rho\) commutes with its
complement, so \(C_M\) is central in the entire
\(U(\mathfrak g)\).

Each \(a_i\) kills the trivial module, hence so
does \(C_M\). On \(M\),
\[
 \operatorname{tr}_M C_M
 =\sum_i B_M(a_i,a^i)=\dim\mathfrak h\ne0.
 \tag{3.4}\label{eq:casimir-trace}
\]
Thus its action is a nonzero module endomorphism.
Its kernel and image are submodules of the simple
module \(M\), so it is invertible. Over an
algebraically closed field Schur's lemma makes
it scalar, and \eqref{eq:casimir-trace} gives the
value. Over a general field invertibility, rather
than a scalar formula, is all that is asserted.
\end{proof}

\begin{theorem}[Casimir contraction of nontrivial simple coefficients]
\label{res:C03}
For a nontrivial finite-dimensional simple
\(\mathfrak g\)-module \(M\),
\[
 H^q(\mathfrak g,M)=0\qquad(q\geq0).
\]
The vanishing is induced by a nullhomotopy for the
central Casimir action on a free resolution of the
trivial module.
\end{theorem}

\begin{proof}
Use the Chevalley--Eilenberg free resolution
\[
 P_q=U(\mathfrak g)\otimes\bigwedge^q\mathfrak g
       \longrightarrow k.
 \tag{3.5}\label{eq:ce-resolution}
\]
Its differential, freeness, and exactness are
proved in \cite[Chapter 6, A06--A07]{homological}:
the PBW filtration identifies the associated graded
resolution with the Koszul resolution, and the
finite-degree lifting argument proves exactness.
Thus its use here has a previously established
proof dependency.

Multiplication by \(C_M\) is a \(U(\mathfrak g)\)-linear
chain endomorphism of \(P_\bullet\), because it is
central; it induces zero on \(k\).
We give the comparison homotopy explicitly at the
lifting level. Put \(h_{-1}=0\). Suppose \(h_i\)
has been chosen below degree \(q\) so that
\[
 C_M=d h_i+h_{i-1}d
\]
there. The map
\(C_M-h_{q-1}d:P_q\to P_q\) has image in
\(\ker d\): applying \(d\), using the already
constructed equation in degree \(q-1\), and using
\(d^2=0\), gives zero. At degree zero the same
statement uses the augmentation. Exactness of
\eqref{eq:ce-resolution} and projectivity of the
free \(P_q\) lift this map through
\(d:P_{q+1}\to\ker d\). Choose that lift as \(h_q\).
Induction constructs a homotopy in all degrees.

Apply \(\operatorname{Hom}_{U(\mathfrak g)}(-,M)\).
For a cochain \(f:P_q\to M\), define
\(Hf=f h_{q-1}\).
Then
\[
 (\delta H+H\delta)f=f C_M=C_M f,
\]
the final equality following from \(U\)-linearity.
Thus the cochain map induced by the invertible
endomorphism \(C_M:M\to M\) is nullhomotopic.
It is an isomorphism on cohomology and also zero
on cohomology; hence every cohomology group is zero.
This proof is the central-action argument of
\cite[Theorem 7.8.9]{weibel}, with its comparison
homotopy made explicit.
\end{proof}

\begin{theorem}[Whitehead's first lemma]\label{res:C04}
For every finite-dimensional \(\mathfrak g\)-module \(M\),
\[
 H^1(\mathfrak g,M)=0.
 \tag{3.6}\label{eq:whitehead-first}
\]
Equivalently, every linear map \(a:\mathfrak g\to M\)
satisfying
\[
 a([x,y])=x\,a(y)-y\,a(x)
\]
has the form \(a(x)=xm\) for some \(m\in M\).
\end{theorem}

\begin{proof}
The low-degree Chevalley--Eilenberg formulas identify
one-cocycles and one-coboundaries exactly as in the
statement; see \cite[Chapter 6, A08]{homological}.
If \(M\) is simple and nontrivial, apply
Theorem~\ref{res:C03}. A simple trivial module is \(k\),
and
\[
 H^1(\mathfrak g,k)=
 \operatorname{Hom}_k(\mathfrak g/[\mathfrak g,\mathfrak g],k)=0
\]
by perfectness from Theorem~\ref{res:B06}.
For arbitrary \(M\), induct along a composition
series. For \(0\to U\to M\to Q\to0\), the cochain
complexes form a degreewise short exact sequence,
since \(\operatorname{Hom}_k(\bigwedge^q\mathfrak g,-)\)
is exact over a field. The cohomology exact sequence
contains
\(H^1(\mathfrak g,U)\to H^1(\mathfrak g,M)
\to H^1(\mathfrak g,Q)\).
Both outer terms vanish by induction, proving the
assertion.
\end{proof}

\begin{theorem}[Weyl complete reducibility]\label{res:C05}
Every submodule of a finite-dimensional
\(\mathfrak g\)-module has a \(\mathfrak g\)-invariant
complement. Consequently every such module is a
finite direct sum of simple modules.
\end{theorem}

\begin{proof}
Let \(0\to W\to V\xrightarrow{p}Q\to0\), and choose
a \(k\)-linear section \(s:Q\to V\).
Its equivariance defect is
\[
 a(x)=\rho_V(x)s-s\rho_Q(x)\in\operatorname{Hom}_k(Q,W),
 \tag{3.7}\label{eq:splitting-cocycle}
\]
since applying \(p\) gives zero. The Hom module has
the action \(x\cdot t=\rho_W(x)t-t\rho_Q(x)\).
Expanding \eqref{eq:splitting-cocycle} at \([x,y]\)
and using the representation identities gives
\(a([x,y])=x\cdot a(y)-y\cdot a(x)\);
the two mixed products cancel.
Whitehead's first lemma supplies \(t:Q\to W\) with
\(a(x)=x\cdot t\).
The adjusted section \(s-t\) has zero defect and
still satisfies \(p(s-t)=1\). Its image is the
desired invariant complement. Split off a simple
submodule, and repeat on the smaller complement;
finite dimension terminates the process.
\end{proof}

This proof distinguishes two statements often both
called semisimplicity: Theorem~\ref{res:B06} concerns
simple ideals of the algebra, whereas
Theorem~\ref{res:C05} concerns arbitrary modules over
it. Nor does the theorem apply to an arbitrary
reductive algebra with no condition on the action of
its center. A central generator can act by a
nonzero nilpotent Jordan block.

\section{Rank-one consequences and functorial Jordan decomposition}

\begin{theorem}[Clebsch--Gordan, duals, and invariant forms]
\label{res:C06}
Over an algebraically closed characteristic-zero field,
every finite-dimensional \(\mathfrak{sl}_2\)-module is
a direct sum of the \(L_n\), and
\[
 L_m\otimes L_n\simeq
 \bigoplus_{j=0}^{\min(m,n)}L_{m+n-2j}.
 \tag{3.8}\label{eq:clebsch-gordan}
\]
Each \(L_n\) is self-dual and has, up to a scalar,
a unique nonzero invariant bilinear form. It is
nondegenerate, symmetric for even \(n\), and
alternating for odd \(n\).
\end{theorem}

\begin{proof}
Simplicity of \(\mathfrak{sl}_2\), the Killing
criterion, and Weyl's theorem give complete
reducibility; Theorem~\ref{res:C01} identifies its
simple summands.
Record the \(h\)-weights by
\(\chi_n(t)=t^n+t^{n-2}+\cdots+t^{-n}\).
The coefficient of \(t^{m+n-2s}\) in
\(\chi_m\chi_n\) counts pairs
\((a,b)\) with \(0\leq a\leq m,0\leq b\leq n,a+b=s\).
If \(m\leq n\), this number is
\[
 \min(m,s)-\max(0,s-n)+1
\]
when the interval is nonempty, and zero otherwise.
The same count results from
\(\sum_{j=0}^{m}\chi_{m+n-2j}\): its \(j\)-th
term contributes exactly when
\(j\leq s\leq m+n-j\). Hence
\(\chi_m\chi_n=\sum_{j=0}^{\min(m,n)}
\chi_{m+n-2j}\).
The characters \(\chi_r\) are linearly independent,
because a nontrivial finite linear relation has a
largest exponent with a nonzero coefficient.
Complete reducibility therefore gives
\eqref{eq:clebsch-gordan}.

In the basis \eqref{eq:sl2-basis}, define
\[
 B(v_i,v_j)=
 \begin{cases}(-1)^i,&i+j=n,\\0,&i+j\ne n.\end{cases}
 \tag{3.9}\label{eq:sl2-form}
\]
For \(h\), invariance
\(B(hv_i,v_j)+B(v_i,hv_j)=0\) follows from the
weights. For \(f\), its only possible nonzero
case is \(i+j=n-1\), and the two terms are
\((-1)^{i+1}+(-1)^i=0\).
For \(e\), its only possible case is \(i+j=n+1\);
the coefficients \(i(n-i+1)\) and
\(j(n-j+1)\) are then equal, and the signs
are opposite. Thus \(B\) is invariant.
Its antidiagonal entries are nonzero, so it
identifies \(L_n\) with \(L_n^*\).
Any invariant form must be antidiagonal by \(h\),
and the \(f\)-equations force successive entries
to change sign. Thus it is a scalar multiple
of \eqref{eq:sl2-form}.
Finally
\(B(v_j,v_i)=(-1)^n B(v_i,v_j)\).
For odd \(n\) its diagonal is zero, giving the
alternating assertion.
\end{proof}

\begin{theorem}[Representations preserve Jordan decomposition]
\label{res:C07}
Let \(k\) have characteristic zero and let
\(\rho:\mathfrak g\to\operatorname{End}_k(V)\)
be finite dimensional, with \(\mathfrak g\) semisimple.
For the abstract Jordan decomposition of
Theorem~\ref{res:B07},
\[
 \rho(x_s)=\rho(x)_s,\qquad
 \rho(x_n)=\rho(x)_n.
 \tag{3.10}\label{eq:representation-jordan}
\]
In particular, each toral subalgebra of \(\mathfrak g\) acts semisimply on \(V\).
Over an algebraically closed field, the restriction of \(\rho\) to such a
subalgebra gives a decomposition of \(V\) into simultaneous weight spaces.
\end{theorem}

\begin{proof}
First take an algebraically closed field and a
simple nonzero \(V\), and put
\(\mathfrak a=\rho(\mathfrak g)\).
It is semisimple. Its normalizer in
\(\operatorname{End}(V)\) acts by derivations on
\(\mathfrak a\). Since all such derivations are
inner, every normalizing operator is \(z+c\),
with \(z\in\mathfrak a\) and \(c\) centralizing
\(\mathfrak a\). Schur's lemma makes \(c\)
a scalar.

Let \(y\in\mathfrak a\). Its matrix semisimple
part \(y_s\) normalizes \(\mathfrak a\):
\(\operatorname{ad}y\) preserves \(\mathfrak a\),
so its polynomial semisimple part does also,
and that part is \(\operatorname{ad}y_s\), by
the commuting semisimple/nilpotent calculation
used in the trace criterion. Write \(y_s=z+c\)
as above. Since \(\mathfrak a\) is perfect,
\(\operatorname{tr}y=\operatorname{tr}z=0\).
Nilpotents have zero trace, so
\(\operatorname{tr}y_s=\operatorname{tr}y=0\).
It follows that \((\dim V)c=0\), hence \(c=0\).
Thus \(y_s\) and \(y_n=y-y_s\) actually belong
to \(\mathfrak a\). Their adjoint operators
on \(\mathfrak a\) are semisimple and nilpotent,
respectively: on the ambient endomorphism
space they have these properties, and
restriction to an invariant subspace preserves
them. They are therefore the abstract Jordan
parts in \(\mathfrak a\).

The kernel of \(\rho\) is a direct ideal
summand of \(\mathfrak g\), by
Theorem~\ref{res:B06}. The adjoint action is
block diagonal for this splitting, so abstract
Jordan decomposition commutes with projection
onto the complementary summand. The preceding
paragraph proves \eqref{eq:representation-jordan}
on \(V\).
For an arbitrary \(V\), Weyl's theorem writes
it as a sum of simple modules. Jordan parts
of block diagonal matrices are the sums of
the parts on their blocks, proving the same
formula.

For a nonclosed field, extend to an algebraic
closure. The two sides of
\eqref{eq:representation-jordan} are scalar
extensions of their \(k\)-linear operators,
by the polynomial and uniqueness construction
of Jordan decomposition. Their equality
therefore descends.

Finally a toral algebra is abelian with
semisimple adjoint action for each element.
The formula just proved makes its representing
operators semisimple. Commuting semisimple
operators over an algebraically closed field
are simultaneously diagonalizable: diagonalize
one generator, observe that its eigenspaces
are stable under all the others, and repeat
inside those eigenspaces. The simultaneous
eigenvalues are linear functionals on the
toral algebra and give its weight spaces.
\end{proof}
